\documentclass{article}

\usepackage{arxiv}
\usepackage[utf8]{inputenc} % allow utf-8 input
\usepackage[T1]{fontenc}    % use 8-bit T1 fonts
\usepackage{hyperref}       % hyperlinks
\usepackage{url}            % simple URL typesetting
\usepackage{booktabs}       % professional-quality tables
\usepackage{amsfonts}       % blackboard math symbols
\usepackage{nicefrac}       % compact symbols for 1/2, etc.
\usepackage{microtype}      % microtypography
\usepackage{lipsum}			% Can be removed after putting your text content
\usepackage{graphicx}
\usepackage{natbib}
\usepackage{doi}

\title{Parallel Matrix Computations in Quantitative Finance}

%\date{November 16, 2025}	% Here you can change the date presented in the paper title

\author{ 
	\href{https://orcid.org/0000-0000-0000-0000}
	{
		\includegraphics[scale=0.06]{orcid.pdf}
		\hspace{1mm}Abigail M. ~Johnson}
		\thanks{} \\
	    Chandra Department of Electrical and Computer Engineering\\
	    The University of Texas at Austin\\
	    Austin, TX 78712 \\
	    \texttt{abigailjohnson@utexas.edu} \\
    }

% Uncomment to remove the date
%\date{}

% Uncomment to override  the `A preprint' in the header
%\renewcommand{\headeright}{Technical Report}
%\renewcommand{\undertitle}{Technical Report}
\renewcommand{\shorttitle}{Parallel Matrix Computations in Quantitative Finance}

%%% Add PDF metadata to help others organize their library
%%% Once the PDF is generated, you can check the metadata with
%%% $ pdfinfo template.pdf
\hypersetup{
pdftitle={A template for the arxiv style},
pdfsubject={q-bio.NC, q-bio.QM},
pdfauthor={Abigail M.~Johnson},
pdfkeywords={First keyword, Second keyword, More},
}

\begin{document}
\maketitle

\begin{abstract}
This paper explores fine, course, and hybrid forms of parallelism in matrix computation workloads common in the quantitative finance industry. By evaluating various parallelism strategies, we observe where concurrency delivers real gains and where contention starts to erode. The benchmark suite is implemented in Python3 and leverages SciPy to make use of the underlying BLAS kernel. The high-level design of the benchmark suite consists of two performance test modules, a command line interface to control them, and a bash script to configure and execute a series of variable performance tests. Workloads leveraging fine grained parallelism saturate quickly while course grained workloads continue to scale with additional CPU cores. The results underpin patterns observed in industry, where the biggest performance gains are made by tuning concurrency to workload structure and available hardware.
\end{abstract}

% keywords can be removed
\keywords{Parallel Computing \and Matrix Operations \and Quantitative Finance
    \and Performance \and Benchmark}

\section{Introduction}
This paper explores fine, course, and hybrid forms of parallelism in matrix computations common in the quantitative finance industry. Matrix operations are a well-studied and highly optimized area, so instead of trying to outperform existing libraries, I focus on how they are employed and scaled in real-life applications. In practice, there are many quantitative finance applications that generalize to two popular workflows: 1) using a single matrix factorization to solve many right hand side (RHS) equations and 2) factorizing and solving many independent matrices. By evaluating variable parallelism strategies within these workload types, we observe where concurrency delivers real gains and where contention starts to erode, and the practical tradeoffs therein.

\section{Approach and Methodology}
\label{sec:headings}
In the design of the benchmark suite, I followed an ethos of simulating real life applications, while using abstraction and simplification where it was reasonable and practical to do so. I designed and implemented two modules to benchmark the effects of variable parallelism strategies in matrix computation workloads commonly employed in quantitative finance.

\subsection{Design}
The high-level design of the benchmark suite consists of two performance test modules, a command line interface to control them, and a bash script to configure and execute a series of variable performance tests. The first module, sameA, models fine-grained parallelism where a single matrix (sameA) is factored once and reused to solve multiple right hand side (RHS) equations; this mirrors common workflows such as Greek calculation and portfolio risk attribution [1][2]. The second module, manyA, simulates course-grained, task-level parallelism and is responsible for solving multiple independent matrices, emulating many-market scenarios such as Monte Carlo engines [1]. To run the benchmark suite, I opted for a configurable command line interface to instantiate both individual performance tests and the collective benchmark suite. This allowed for ease of implementation while promoting flexibility and future extensibility. 


\subsection{Implementation}
The benchmark suite is implemented in Python3 and leverages SciPy to make use of the underlying BLAS kernel for matrix operations. I opted for Python to focus on workload emulation instead of low-level or algorithmic optimizations via Cuda or OpenMP. This was a notable tradeoff, but a deliberate choice, as the chosen stack better emulates the development environment of a typical quantitative researcher where vendor-optimized kernels such as BLAS and LAPACK are the standard backend <[ref]>. 
Both modules, sameA and manyA, leverage the “spd” helper to construct random symmetric positive definite (SPD) matrix(s) of size NxN using the formula  A =X  TX +  I. The sameA module constructs a single SPD matrix A, factorizes it once, and solves A x = b for ‘S’ right hand sides using ‘inner’ BLAS kernel threads. The manyA module constructs and factorizes ‘S’ SPD matrices and solves  As xs = b where eb is a vector of ones. Parallelism in manyA comes from using ProcessPoolExecutor to distribute tasks across ‘outer’ worker processes and assigning ‘inner’ per process BLAS kernel threads. Having both ‘inner’ and ‘outer’ configurable parameters enables the manyA module to express both course and hybrid parallelism. 

\section{Results}
\label{sec:headings}
Table 1 below details the results of running the benchmark suite configured for 256x256 matrix and 256 scenarios on two machines: a 2020 Apple M1 with eight CPU Cores and a 2024 Apple M4 Pro with twelve CPU Cores. As illustrated in <Figure Y>, performance improvements for fine grained parallelism quickly plateaus, and we see near identical scaling patterns on both machines, despite core count differences. In these runs, there is a slight performance boost increasing from one to two BLAS threads, but it degrades thereafter. Course grained parallelism scaled up to four processes on the M1, and degraded thereafter. On the M4 Pro, it scaled almost linearly up to eight. For course grained runs, throughput roughly doubles from the M1 to the M4 Pro. Hybrid parallelism provides modest gains over coarse grained parallelism on the M1 machine, but comparatively underperforms on the M4 Pro.

\section{Discussion}
\label{sec:headings}
The results underpin patterns observed in industry where the biggest performance gains are made by tuning concurrency to workload structure and available hardware. The results of the fine-grained benchmarks are unsurprising, and highlight that after initial factorization, the remaining computation is dominated by memory-bound tasks rather than CPU-bound tasks. This demonstrates that the BLAS kernel is almost entirely bandwidth limited, and additional threads only create more contention. The results of course-grained benchmarks are more pronounced. On both machines, scenario-parallel computations scale well with core count. This is especially pronounced on the M4 Pro machine, and highlights architecture improvements introduced in the generation such as higher memory bandwidth and larger shared caches <ref>. The results from the hybrid benchmarks are more nuanced, and appear to highlight differences in architecture and where memory pressure first becomes limiting.


\section{Conclusions and Future Work}
\label{sec:headings}
Overall, the project and benchmark suite was successfully scoped, designed, and implemented to address the original problem statement: evaluating the effects of varying parallelism strategies in matrix computation workloads commonly employed in the quantitative finance industry. The high-level design of the benchmark suite consists of two performance test modules, a command line interface to control them, and a bash script to configure and execute a series of variable performance tests. Running the benchmark suite on two machines, a 2020 Apple M1 with eight CPU cores and a 2024 Apple M4 Pro with twelve CPU cores, demonstrate that workloads leveraging fine grained parallelism saturate quickly, while course grained workloads continue to scale with additional CPU cores. The results underpin patterns observed in industry where the biggest performance gains are made by tuning concurrency to workload structure and available hardware. While time did not permit their inclusion in the initial scope of this project, in the future, there are a number of topics I would like to explore further: comparing the benchmark with the Rust implemented equivalent, latest Python equivalent (currently unstable release), optimization with cuPy, and further investigation into the nuance of the hybrid case. 



\section{Examples of citations, figures, tables, references}
\label{sec:others}

\subsection{Citations}
Citations use \verb+natbib+. The documentation may be found at
\begin{center}
	\url{http://mirrors.ctan.org/macros/latex/contrib/natbib/natnotes.pdf}
\end{center}

Here is an example usage of the two main commands (\verb+citet+ and \verb+citep+): Some people thought a thing \citep{kour2014real, hadash2018estimate} but other people thought something else \citep{kour2014fast}. Many people have speculated that if we knew exactly why \citet{kour2014fast} thought this\dots

\subsection{Figures}
See Figure \ref{fig:fig1}. Here is how you add footnotes. \footnote{Sample of the first footnote.}

\begin{figure}
	\centering
	\fbox{\rule[-.5cm]{4cm}{4cm} \rule[-.5cm]{4cm}{0cm}}
	\caption{Sample figure caption.}
	\label{fig:fig1}
\end{figure}

\subsection{Tables}
See awesome Table~\ref{tab:table}.

The documentation for \verb+booktabs+ (`Publication quality tables in LaTeX') is available from:
\begin{center}
	\url{https://www.ctan.org/pkg/booktabs}
\end{center}


\begin{table}
	\caption{Sample table title}
	\centering
	\begin{tabular}{lll}
		\toprule
		\multicolumn{2}{c}{Part}                   \\
		\cmidrule(r){1-2}
		Name     & Description     & Size ($\mu$m) \\
		\midrule
		Dendrite & Input terminal  & $\sim$100     \\
		Axon     & Output terminal & $\sim$10      \\
		Soma     & Cell body       & up to $10^6$  \\
		\bottomrule
	\end{tabular}
	\label{tab:table}
\end{table}


\bibliographystyle{unsrtnat}
\bibliography{references}  %%% Uncomment this line and comment out the ``thebibliography'' section below to use the external .bib file (using bibtex) .


%%% Uncomment this section and comment out the \bibliography{references} line above to use inline references.
% \begin{thebibliography}{1}

% 	\bibitem{kour2014real}
% 	George Kour and Raid Saabne.
% 	\newblock Real-time segmentation of on-line handwritten arabic script.
% 	\newblock In {\em Frontiers in Handwriting Recognition (ICFHR), 2014 14th
% 			International Conference on}, pages 417--422. IEEE, 2014.

% 	\bibitem{kour2014fast}
% 	George Kour and Raid Saabne.
% 	\newblock Fast classification of handwritten on-line arabic characters.
% 	\newblock In {\em Soft Computing and Pattern Recognition (SoCPaR), 2014 6th
% 			International Conference of}, pages 312--318. IEEE, 2014.

% 	\bibitem{hadash2018estimate}
% 	Guy Hadash, Einat Kermany, Boaz Carmeli, Ofer Lavi, George Kour, and Alon
% 	Jacovi.
% 	\newblock Estimate and replace: A novel approach to integrating deep neural
% 	networks with existing applications.
% 	\newblock {\em arXiv preprint arXiv:1804.09028}, 2018.

% \end{thebibliography}


\end{document}
